\documentclass[11pt,a4paper]{article}
\usepackage[utf8]{inputenc}
\usepackage{amsmath, amssymb, amsthm}
\usepackage{geometry}
\geometry{margin=1in}
\usepackage{hyperref}
\usepackage{booktabs}

\title{DRESS and the WL Hierarchy: Climbing One Deletion at a Time}
\author{Eduar Castrillo Velilla \\ \texttt{eduarcastrillo@gmail.com}}
\date{}

\newtheorem{theorem}{Theorem}
\newtheorem{proposition}[theorem]{Proposition}
\newtheorem{corollary}[theorem]{Corollary}
\newtheorem{definition}[theorem]{Definition}

\begin{document}

\maketitle

\begin{abstract}
The Cai--F\"urer--Immerman (CFI) construction provides the canonical family of hard instances for the Weisfeiler--Leman (WL) hierarchy: distinguishing the two non-isomorphic CFI graphs over a base graph $G$ requires $k$-WL where $k$ meets or exceeds the treewidth of $G$. In this paper, we introduce $\Delta^\ell$-DRESS, which applies $\ell$ levels of iterated node deletion to the DRESS continuous structural refinement framework. $\Delta^\ell$-DRESS runs Original-DRESS on all $\binom{n}{\ell}$ subgraphs obtained by removing $\ell$ nodes, and compares the resulting histograms. We show empirically on the canonical CFI benchmark family that Original-DRESS ($\Delta^0$) already distinguishes $\text{CFI}(K_3)$ (requiring 2-WL), and that each additional deletion level extends the range by one WL level: $\Delta^1$ reaches 3-WL, $\Delta^2$ reaches 4-WL, and $\Delta^3$ reaches 5-WL, distinguishing CFI pairs over $K_n$ for $n = 3, \ldots, 6$. Crucially, $\Delta^3$ fails on $\text{CFI}(K_7)$ (requiring 6-WL), confirming a sharp boundary at $(\ell+2)$-WL. The computational cost is $\mathcal{O}\bigl(\binom{n}{\ell} \cdot I \cdot m \cdot d_{\max}\bigr)$---polynomial in $n$ for fixed $\ell$. These results establish $\Delta^\ell$-DRESS as a practical framework for systematically climbing the WL hierarchy on the canonical CFI benchmark family.
\end{abstract}

\section{Introduction}

The Weisfeiler--Leman (WL) hierarchy \cite{weisfeiler1968reduction} is the standard yardstick for measuring the expressiveness of graph isomorphism algorithms. The $k$-WL test operates on $k$-tuples of vertices and runs in $\mathcal{O}(n^{k+1})$ time per iteration \cite{cai1992optimal}. While powerful in theory, even 3-WL at $\mathcal{O}(n^4)$ is impractical for large graphs.

The CFI construction \cite{cai1992optimal} provides the canonical hard instances: for any base graph $G$, it produces a pair of non-isomorphic graphs $\text{CFI}(G)$ and $\text{CFI}'(G)$ that $k$-WL cannot distinguish whenever $k$ is below the treewidth of $G$. In particular, $\text{CFI}(K_n)$ requires $(n{-}1)$-WL, making these the hardest known instances for the WL hierarchy.

In the companion paper \cite{castrillo2025dress}, we introduced the DRESS family of continuous structural refinement algorithms and showed that $\Delta$-DRESS---which runs DRESS on each single-node-deleted subgraph---empirically distinguishes strongly regular graphs that confound 3-WL. The connection to the Kelly--Ulam reconstruction conjecture \cite{kelly1942congruence, ulam1960collection} was noted: the multiset of node-deleted DRESS fingerprints is a continuous relaxation of the reconstruction deck.

In this paper, we generalize $\Delta$-DRESS to $\Delta^\ell$-DRESS by applying $\ell$ levels of iterated node deletion. We show that each deletion level adds approximately one WL level of discriminative power on the CFI family: $\Delta^0$ (Original-DRESS) already distinguishes $\text{CFI}(K_3)$, $\Delta^1$ reaches $\text{CFI}(K_4)$, $\Delta^2$ reaches $\text{CFI}(K_5)$, and $\Delta^3$ reaches $\text{CFI}(K_6)$.

\subsection{Contributions}
\begin{enumerate}
    \item We define $\Delta^\ell$-DRESS formally and analyze its complexity.
    \item We show that $\Delta^0$ (Original-DRESS) already distinguishes $\text{CFI}(K_3)$---a result not previously reported.
    \item We present empirical results showing that $\Delta^\ell$-DRESS distinguishes $\text{CFI}(K_{\ell+3})$ for each $\ell = 0, 1, 2, 3$, and fails at the next level, covering WL levels 2 through 5.
    \item We observe that each deletion level adds approximately one WL level of expressiveness on the CFI family.
    \item We discuss the theoretical implications and open questions raised by these results.
\end{enumerate}

\section{Preliminaries}

\subsection{DRESS}

The Original-DRESS equation \cite{castrillo2018dynamicstructuralsimilaritygraphs} is a parameter-free, non-linear dynamical system on edges:

\begin{equation}
d_{uv}^{(t+1)} = \frac{\sum_{x \in N[u] \cap N[v]} \bigl(d_{ux}^{(t)} + d_{xv}^{(t)}\bigr)}{\|u\|^{(t)} \cdot \|v\|^{(t)}}
\end{equation}

where $\|u\|^{(t)} = \sqrt{\sum_{x \in N[u]} d_{ux}^{(t)}}$ is the node norm and $N[u] = N(u) \cup \{u\}$ is the closed neighborhood. Self-loops are added to every node; the invariant $d_{uu} = 2$ holds at every iteration. The equation converges to a unique fixed point $d^* \in [0, 2]^{|E|}$ for any positive initialization \cite{castrillo2025dress}.

The graph fingerprint can be represented equivalently as the sorted vector $\text{sort}(d^*)$ or as the histogram $h(d^*)$; both uniquely identify the multiset of converged edge values. Since DRESS values are bounded in $[0, 2]$ and convergence tolerance is $\epsilon = 10^{-6}$, each edge value maps to one of $\lceil 2/\epsilon \rceil = 2 \times 10^{6}$ integer bins (the bin count is determined by the convergence tolerance), so the histogram is a compact bin-count vector. In this work we use the histogram representation. Two graphs are declared non-isomorphic if their fingerprints differ.

\subsection{\texorpdfstring{$\Delta$}{Delta}-DRESS}

$\Delta$-DRESS \cite{castrillo2025dress} runs DRESS on each node-deleted subgraph:

\begin{definition}[$\Delta$-DRESS]
Given a graph $G = (V, E)$ and a DRESS variant $\mathcal{F}$, the $\Delta$-DRESS fingerprint is the histogram obtained by running $\mathcal{F}$ on each node-deleted subgraph and accumulating every edge similarity value into a single bin-count vector:
\[
\Delta\text{-DRESS}(\mathcal{F}, G) = h\!\left(\bigsqcup_{v \in V} \mathcal{F}(G \setminus \{v\})\right)
\]
where $G \setminus \{v\}$ is the subgraph induced by $V \setminus \{v\}$, $\bigsqcup$ denotes multiset union, and $h$ maps each value to its integer bin in a fixed-size array.
\end{definition}

\subsection{CFI Graphs}

The CFI construction \cite{cai1992optimal} is defined over a base graph $G = (V_G, E_G)$. For each vertex $v \in V_G$, a ``gadget'' replaces $v$ with a set of vertices encoding its incident edges. Two non-isomorphic graphs $\text{CFI}(G)$ and $\text{CFI}'(G)$ are produced by twisting the parity of one gadget. The key property is:

\begin{theorem}[Cai, F\"urer, Immerman \cite{cai1992optimal}]
The $k$-WL algorithm cannot distinguish $\text{CFI}(G)$ from $\text{CFI}'(G)$ whenever $k < \text{tw}(G)$, where $\text{tw}(G)$ is the treewidth of $G$.
\end{theorem}

For the complete graph $K_n$, the treewidth is $n - 1$, so $\text{CFI}(K_n)$ requires $(n{-}1)$-WL. In particular:
\begin{itemize}
    \item $\text{CFI}(K_3)$: requires 2-WL.
    \item $\text{CFI}(K_5)$: requires 4-WL, operating on 4-tuples---$\mathcal{O}(n^5)$ per iteration.
    \item $\text{CFI}(K_6)$: requires 5-WL, operating on 5-tuples---$\mathcal{O}(n^6)$ per iteration.
\end{itemize}

\section{\texorpdfstring{$\Delta^\ell$}{Delta-ell}-DRESS}

We now define the iterated deletion operator.

\begin{definition}[$\Delta^\ell$-DRESS]
For a graph $G = (V, E)$, a DRESS variant $\mathcal{F}$, and a deletion depth $\ell \ge 0$:
\[
\Delta^\ell\text{-DRESS}(\mathcal{F}, G) = h\!\left(\bigsqcup_{\substack{S \subset V,\; |S|=\ell}} \mathcal{F}(G \setminus S)\right)
\]
where $G \setminus S$ is the subgraph induced by $V \setminus S$, and $h$ maps the pooled edge values into a histogram.
\end{definition}

At depth $\ell$, $\Delta^\ell$-DRESS computes DRESS on all $\binom{n}{\ell}$ subgraphs obtained by removing exactly $\ell$ vertices. Every individual edge similarity value from every subgraph is added to a single histogram of $G$. Two graphs are compared by checking equality of their histograms. Individual subgraph fingerprints are not recoverable from the pooled histogram. The memory footprint is constant regardless of $\ell$, whereas storing the raw multiset of $\binom{n}{\ell} \cdot |E|$ floating-point values would be prohibitive. It is noteworthy that this fixed-size representation, combined with the deletion operator $\Delta^\ell$, empirically matches the discriminative power of $k$-WL methods that require $\mathcal{O}(n^{k+1})$ storage.

\begin{proposition}[Complexity]
The time complexity of $\Delta^\ell$-DRESS is
\[
\mathcal{O}\!\left(\binom{n}{\ell} \cdot I \cdot m \cdot d_{\max}\right)
\]
where $n = |V|$, $m = |E|$, $d_{\max}$ is the maximum degree, and $I$ is the number of DRESS iterations per subgraph. The $\binom{n}{\ell}$ subgraph computations are entirely independent and embarrassingly parallel.
\end{proposition}

\paragraph{Comparison with \texorpdfstring{$k$}{k}-WL.} The cost of $k$-WL is $\mathcal{O}(n^{k+1})$ per iteration. For $\text{CFI}(K_6)$, 5-WL costs $\mathcal{O}(n^6)$, while $\Delta^3$-DRESS costs $\mathcal{O}(n^3 \cdot I \cdot m \cdot d_{\max}) = \mathcal{O}(n^3 \cdot m \cdot d_{\max})$ (with $I$ constant).

\section{Empirical Results}

All experiments use Original-DRESS as the base variant $\mathcal{F}$, with convergence tolerance $\epsilon = 10^{-6}$ and a maximum of 100 iterations. In every case, convergence was reached before the iteration limit.

\subsection{\texorpdfstring{$\Delta^\ell$}{Delta-ell}-DRESS on \texorpdfstring{CFI$(K_n)$}{CFI(Kn)}}

We test $\Delta^\ell$-DRESS for $\ell = 0, 1, 2, 3$ on CFI pairs over complete graphs $K_n$ ($n = 3, \ldots, 10$). Higher deletion depths on larger base graphs were not executed due to time constraints (marked --- in the table).

\begin{table}[h]
\centering
\caption{$\Delta^\ell$-DRESS results on CFI graph pairs. $\checkmark$ = pair distinguished; $\times$ = pair not distinguished; --- = not executed due to time constraints. The WL requirement column shows the minimum $k$-WL level needed to distinguish the pair.}
\label{tab:cfi-results}
\begin{tabular}{@{}lcccccc@{}}
\toprule
Base graph & $|V_{\text{CFI}}|$ & WL req. & $\Delta^0$ & $\Delta^1$ & $\Delta^2$ & $\Delta^3$ \\
\midrule
$K_3$   & 6     & 2-WL & $\checkmark$  & $\checkmark$  & $\checkmark$  & $\checkmark$  \\
$K_4$   & 16    & 3-WL & $\times$      & $\checkmark$  & $\checkmark$  & $\checkmark$  \\
$K_5$   & 40    & 4-WL & $\times$      & $\times$      & $\checkmark$  & $\checkmark$  \\
$K_6$   & 96    & 5-WL & $\times$      & $\times$      & $\times$      & $\checkmark$  \\
$K_7$   & 224   & 6-WL & $\times$      & $\times$      & $\times$      & $\times$      \\
$K_8$   & 512   & 7-WL & $\times$      & $\times$      & ---            & ---            \\
$K_9$   & 1152  & 8-WL & $\times$      & $\times$      & ---            & ---            \\
$K_{10}$& 2560  & 9-WL & $\times$      & $\times$      & ---            & ---            \\
\bottomrule
\end{tabular}
\end{table}

\subsection{Analysis}

The results in Table~\ref{tab:cfi-results} reveal a clear staircase pattern:

\begin{enumerate}
    \item \textbf{$\Delta^0$ (Original-DRESS)} distinguishes $\text{CFI}(K_3)$, which requires 2-WL. This is a notable result: Original-DRESS---with no node deletion---already exceeds 1-WL on CFI instances, consistent with the prism vs.\ $K_{3,3}$ result in \cite{castrillo2025dress}. It fails on all larger CFI pairs.

    \item \textbf{$\Delta^1$} extends the range to $\text{CFI}(K_4)$ (3-WL). This is consistent with the results in \cite{castrillo2025dress}, where $\Delta$-DRESS was shown to distinguish SRG pairs that confound 3-WL. It fails on $\text{CFI}(K_5)$.

    \item \textbf{$\Delta^2$} reaches $\text{CFI}(K_5)$ (4-WL).

    \item \textbf{$\Delta^3$} reaches $\text{CFI}(K_6)$ (5-WL) but fails on $\text{CFI}(K_7)$ (6-WL), confirming that the staircase has a sharp boundary at $(\ell+2)$-WL. The $\text{CFI}(K_7)$ test required processing all $\binom{224}{3} = 1{,}848{,}224$ subgraphs ($\approx$19.3 billion edge values).
\end{enumerate}

\paragraph{WL level per deletion depth.} The tested data shows:

\begin{center}
\begin{tabular}{@{}ccc@{}}
\toprule
Deletion depth $\ell$ & Max WL matched & Effective WL \\
\midrule
0 & 2-WL & $\ell + 2$ \\
1 & 3-WL & $\ell + 2$ \\
2 & 4-WL & $\ell + 2$ \\
3 & 5-WL & $\ell + 2$ \\
\bottomrule
\end{tabular}
\end{center}

The pattern is strikingly regular: $\Delta^\ell$-DRESS empirically matches $(\ell+2)$-WL on the CFI family. Each deletion level adds one WL level of expressiveness, with the base DRESS contributing 2-WL on its own. The computational cost grows as $\binom{n}{\ell}$ (polynomial in $n$ for fixed $\ell$), while the equivalent $(\ell+2)$-WL cost grows as $n^{\ell+3}$.

\section{Discussion}

\subsection{Relationship to Existing Work}

\paragraph{Subgraph GNNs.} Methods such as ESAN \cite{bevilacqua2022equivariant} and GNN-AK+ \cite{zhao2022from} also use node-deleted or node-marked subgraphs to boost expressiveness. However, these are supervised methods that learn aggregation functions from data. $\Delta^\ell$-DRESS is entirely unsupervised: the aggregation is the deterministic DRESS fixed point, and the histogram comparison is parameter-free. This makes $\Delta^\ell$-DRESS a canonical baseline for the expressiveness of subgraph-based approaches.

\paragraph{Reconstruction conjecture.} The Kelly--Ulam conjecture \cite{kelly1942congruence, ulam1960collection} states that graphs with $n \ge 3$ are determined by their deck of node-deleted subgraphs. Conceptually, $\Delta^1$-DRESS computes a continuous relaxation of this deck: each ``card'' is the DRESS fingerprint of a node-deleted subgraph. In practice, all edge values are accumulated into a fixed-size histogram for memory efficiency, though individual cards are not recoverable.

\paragraph{Fixed-parameter tractability.} For fixed $\ell$, $\Delta^\ell$-DRESS runs in polynomial time. The empirical pattern $\Delta^\ell \approx (\ell+2)$-WL suggests that the minimum deletion depth needed to distinguish $\text{CFI}(K_n)$ grows as $n - 2$, which is linear in $n$. While each additional deletion level is polynomial, the number of required levels grows with the base graph size.

\subsection{Open Questions}

\begin{enumerate}
    \item \textbf{Does the $\Delta^\ell \approx (\ell+2)$-WL pattern continue?} We have confirmed that $\Delta^3$ fails on $\text{CFI}(K_7)$ (6-WL). Testing $\Delta^4$ on $\text{CFI}(K_7)$ would confirm whether the staircase extends to $\ell = 4$.

    \item \textbf{Is the pattern specific to CFI, or does it hold more broadly?} Other hard graph families may behave differently.

    \item \textbf{Are there graph families where $\Delta^\ell$-DRESS fails for all constant $\ell$?} CFI graphs over graphs of unbounded treewidth would be natural candidates.

    \item \textbf{Does combining $\Delta^\ell$ with other DRESS variants yield further gains?} The DRESS family includes Motif-DRESS and Generalized-DRESS \cite{castrillo2025dress}. Whether a stronger base variant $\mathcal{F}$ shifts the staircase upward or merely overlaps with the deletion mechanism is an open question.
\end{enumerate}

\section{Conclusion}

We introduced $\Delta^\ell$-DRESS, which applies $\ell$ levels of iterated node deletion to the DRESS framework. Original-DRESS ($\Delta^0$) already distinguishes $\text{CFI}(K_3)$, and each additional deletion level adds one WL level: $\Delta^1$ matches 3-WL, $\Delta^2$ matches 4-WL, and $\Delta^3$ matches 5-WL but fails on 6-WL, establishing the empirical pattern $\Delta^\ell \approx (\ell+2)$-WL on the CFI family with a sharp upper boundary. These results establish $\Delta^\ell$-DRESS as a systematic, scalable method for climbing the WL hierarchy on the canonical CFI benchmarks. An open-source implementation is available at \url{https://github.com/velicast/dress-graph}.

\bibliographystyle{plain}
\bibliography{refs}

\end{document}
